\documentclass[11pt,a4paper]{article}

% Fonts: TeX Gyre Termes (Times equivalent) with matching math
\usepackage{fontspec}
\usepackage{unicode-math}
\setmainfont{texgyretermes}[
  Extension      = .otf,
  UprightFont    = *-regular,
  BoldFont       = *-bold,
  ItalicFont     = *-italic,
  BoldItalicFont = *-bolditalic,
]
\setmathfont{texgyretermes-math.otf}
\frenchspacing

% Body size / leading: 11.5pt / 13.5pt
\makeatletter
\renewcommand{\@xpt}{11.5}
\renewcommand{\@xiipt}{13.5}
\renewcommand\normalsize{%
  \@setfontsize\normalsize\@xpt\@xiipt
  \abovedisplayskip 10\p@ \@plus2\p@ \@minus5\p@
  \abovedisplayshortskip \z@ \@plus3\p@
  \belowdisplayshortskip 6\p@ \@plus3\p@ \@minus3\p@
  \belowdisplayskip \abovedisplayskip
  \let\@listi\@listI}
\makeatother
\normalsize

% Page layout
\usepackage[margin=1in]{geometry}

% Math
\usepackage{amsmath,amsthm}

% Tables
\usepackage{array}
\newcolumntype{L}[1]{>{\raggedright\arraybackslash}p{#1}}
\usepackage{booktabs}
\setlength{\extrarowheight}{2pt}

% Float placement
\usepackage{float}

% Captions: "Fig. N." / "Table N." period separator, label smaller
\usepackage{caption}
\captionsetup{figurename=Fig.,labelsep=period,margin=1.2cm,labelfont=footnotesize}

% Lists
\usepackage{enumitem}

% Section heading style (lighter than article default)
\usepackage{titlesec}
\titleformat{\section}[block]{\fontsize{13pt}{15pt}\bfseries}{\thesection.}{0.6em}{}
\titlespacing*{\section}{0pt}{18pt}{8pt}
\titleformat{\subsection}[block]{\fontsize{11.5pt}{13.5pt}\bfseries\itshape}{\thesubsection.}{0.6em}{}
\titlespacing*{\subsection}{0pt}{14pt}{6pt}

% Graphics + TikZ + pgfplots (in-document figures)
\usepackage{graphicx}
\usepackage{tikz}
\usetikzlibrary{positioning,arrows,calc}
\usepackage{pgfplots}
\pgfplotsset{compat=1.18}

% Paragraph rules (indented first line, no gap, indentfirst on)
\usepackage{indentfirst}
\emergencystretch=1em
\widowpenalty=10000
\clubpenalty=10000
\usepackage{needspace}

% Float / display-math spacing
\setlength{\intextsep}{8pt plus 2pt minus 2pt}
\setlength{\floatsep}{8pt plus 2pt minus 2pt}
\setlength{\textfloatsep}{10pt plus 2pt minus 2pt}
\AtBeginDocument{%
  \setlength{\abovedisplayskip}{18pt plus 3pt minus 2pt}%
  \setlength{\belowdisplayskip}{18pt plus 3pt minus 2pt}%
  \setlength{\abovedisplayshortskip}{14pt plus 2pt minus 1pt}%
  \setlength{\belowdisplayshortskip}{14pt plus 2pt minus 1pt}%
}

% Tables: 1.05 row stretch, content at \footnotesize
\renewcommand{\arraystretch}{1.05}
\AtBeginEnvironment{tabular}{\footnotesize}

% Microtype
\usepackage{microtype}

% Citations: superscript numeric, with superscript labels in bibliography
\usepackage[superscript,biblabel]{cite}

% Bibliography: hanging indent with the label sitting in its own column
\makeatletter
\renewenvironment{thebibliography}[1]
  {\section*{\refname}%
   \fontsize{10pt}{12pt}\selectfont
   \list{\@biblabel{\@arabic\c@enumiv}}%
        {\settowidth\labelwidth{\@biblabel{#1}}%
         \leftmargin\labelwidth
         \advance\leftmargin-0.8em
         \advance\leftmargin\labelsep
         \setlength\itemindent{1em}%
         \setlength\itemsep{6pt plus 1pt minus 1pt}%
         \setlength\parsep{0pt}%
         \usecounter{enumiv}%
         \let\p@enumiv\@empty
         \renewcommand\theenumiv{\@arabic\c@enumiv}}%
   \sloppy\clubpenalty4000\widowpenalty4000\sfcode`\.\@m}
  {\def\@noitemerr{\@latex@warning{Empty `thebibliography' environment}}%
   \endlist}
\makeatother

\usepackage{etoolbox}

% Abstract: 10pt/13pt, run-in "Abstract." lead, inset 0.5in each side
\renewenvironment{abstract}{%
  \fontsize{10pt}{13pt}\selectfont
  \begin{list}{}{%
    \setlength{\leftmargin}{0.5in}%
    \setlength{\rightmargin}{0.5in}%
    \setlength{\listparindent}{0pt}%
    \setlength{\parsep}{6pt plus 2pt minus 1pt}%
    \setlength{\topsep}{0pt}%
    \setlength{\partopsep}{0pt}%
  }\item[]\noindent\textbf{Abstract.}\hspace{0.5em}\ignorespaces
}{\end{list}}

% Colours
\usepackage{xcolor}

% Hyperlinks + PDF metadata
\usepackage{hyperref}
\hypersetup{
  colorlinks=true,
  allcolors=blue!50!black,
  pdfauthor={K.R. Ryan},
  pdftitle={The Geometric Siphon II: Directional Properties},
  pdfsubject={Concentrated liquidity, directional and exit asymmetry},
  pdfkeywords={concentrated liquidity, geometric residual, directional asymmetry, exit asymmetry, Uniswap V3, Aerodrome}
}

% Running header (author left, page number centred footer)
\usepackage{fancyhdr}
\pagestyle{fancy}
\fancyhf{}
\fancyhead[L]{\small\itshape K.R.\ Ryan}
\fancyhead[R]{}
\fancyfoot[C]{\thepage}
\renewcommand{\headrulewidth}{0pt}
\fancypagestyle{plain}{%
  \fancyhf{}%
  \fancyfoot[C]{\thepage}%
  \renewcommand{\headrulewidth}{0pt}%
}

% Footnote pinning
\usepackage[bottom]{footmisc}

% Theorem environments
\newtheorem{theorem}{Theorem}
\setcounter{theorem}{3} % Start at Theorem 4

\theoremstyle{definition}
\newtheorem*{notation}{Notation}

% QED symbol (filled black square)
\renewcommand\qedsymbol{\ensuremath{\mdlgblksquare}}

% Code-block environment (cream-tinted, footnotesize)
\usepackage{mdframed}
\newenvironment{codeblock}{%
  \begin{mdframed}[%
    backgroundcolor=brown!3!white,%
    linewidth=0pt,%
    innertopmargin=6pt,%
    innerbottommargin=6pt,%
    innerleftmargin=12pt,%
    innerrightmargin=12pt,%
    leftmargin=0.075\linewidth,%
    rightmargin=0.075\linewidth,%
    skipabove=6pt,%
    skipbelow=6pt%
  ]\footnotesize%
}{%
  \end{mdframed}%
}

% Paragraph spacing: indented first line, no gap between paragraphs
\setlength{\parindent}{18pt}
\setlength{\parskip}{0pt}

\title{The Geometric Siphon II: Directional Properties}

\begin{document}
\sloppy
\thispagestyle{plain}

\begin{center}
{\huge The Geometric Siphon II:\\[6pt]
Directional Properties}

\vspace{14pt}
{\large K.R.\ Ryan}\\[4pt]
{\fontsize{9pt}{11pt}\selectfont Independent Researcher \quad|\quad ORCID - 0009-0004-6295-7040 \quad|\quad \href{mailto:gancor@gancor.xyz}{gancor@gancor.xyz}}\\[4pt]
{\fontsize{9pt}{11pt}\selectfont March 2026}

\vspace{12pt}
{\fontsize{9pt}{11pt}\selectfont
\begin{minipage}{0.72\linewidth}\centering
\textbf{Keywords}\hspace{0.6em}geometric residual \textperiodcentered\ directional asymmetry \textperiodcentered\ exit asymmetry \textperiodcentered\ Uniswap V3 \textperiodcentered\ Aerodrome \textperiodcentered\ DeFi mechanism design\\[4pt]
\textbf{JEL:} G12, G14, G19, G23
\end{minipage}}
\end{center}

\vspace{-1pt}
\centerline{\scalebox{1.5}{$\therefore$}}
\vspace{14pt}

\begin{abstract}
In a companion paper\cite{ryan2026a}, I identified the Geometric Siphon, a mechanism by which concentrated liquidity (CL) positions sharing a depositor-level token balance transfer capital between positions under autonomous rebalancing, with token ratio mismatches between old and new tick ranges producing residuals that flow through the shared balance. This paper establishes three results on the directional properties of that residual. Theorem~4 shows that the absolute residual increases strictly with displacement from the midpoint of the original range, up to the range boundary. Theorems~5 and~6 specialise to pools pairing a volatile asset at token0 with a stablecoin at token1 in Uniswap V3's pair ordering. Under symmetric in-range sqrt-price displacements, the up-move position retains strictly more USD-denominated value than the down-move; under symmetric displacement past the range boundary, a below-range exit retains strictly less USD value than an above-range exit. Residual magnitudes scale with position value, so the USD-denominated absolute residual inherits these asymmetries by a scaling argument under the same domain assumptions. All three results are verified by Foundry tests against a harness implementing exact Uniswap V3 tick-math and amount-equation constants. Across 36{,}337 rebalance events on 58 positions over 18 days on Aerodrome (Base), observed dust flows are consistent with the predicted direction; rising markets show USD-denominated flow about $5\times$ larger than falling markets despite a near-identical frequency of positive dust events. The dataset is unsegmented by pool composition, so these figures support the prediction at production scale without strictly confirming Theorems~5 and~6 in their volatile/stablecoin domain.
\end{abstract}

\vspace{24pt}

\begin{list}{}{%
  \setlength{\leftmargin}{0.5in}%
  \setlength{\rightmargin}{0.5in}%
  \setlength{\listparindent}{0pt}%
  \setlength{\topsep}{0pt}%
  \setlength{\partopsep}{0pt}%
}\item[]\fontsize{10pt}{13pt}\selectfont\noindent\textbf{Note on supersession.} This preprint and its companion \emph{The Geometric Siphon: Emergent Capital Reallocation in Concentrated Liquidity Portfolios} have been consolidated into a single revised manuscript (\emph{The Geometric Siphon: Existence, Equilibrium, and Directional Properties of the Residual in Concentrated Liquidity Portfolios}). The Foundry verification suite and the public dataset now align with the consolidated manuscript.\footnote{Consolidated manuscript: \url{https://github.com/g4nc0r/the-geometric-siphon}. Foundry suite (Code DOI): \url{https://doi.org/10.5281/zenodo.19905208}. Public dataset (Data DOI): \url{https://doi.org/10.5281/zenodo.19904832}.} The empirical findings reported in this preprint remain the citable record for this version of the analysis.\end{list}

\clearpage

\section{Introduction}\label{sec:introduction}

The Geometric Siphon is a mechanism in concentrated liquidity (CL) portfolio management: when multiple CL positions share a depositor-level token balance, autonomous rebalancing produces residuals from token ratio mismatches between old and new tick ranges, and those residuals flow through the shared balance as inter-position capital transfers. Ryan\cite{ryan2026a} established that a CL rebalance produces a zero geometric residual if and only if the old and new ranges demand the same token ratio at current price (Theorem~1), that single-pool positions sharing a contract balance converge to equal value under repeated rebalancing (Theorem~2), and that positions rebalancing without swaps decay geometrically (Theorem~3). Those results characterise the existence, equilibrium and extinction behaviour of the siphon. They do not address how the residual responds to displacement magnitude, price direction or exit type.

Three theorems below address those questions. Theorem~4 (Residual Monotonicity) gives strict monotonicity of $|\Delta R|$ in displacement up to the range boundary; the result holds for any pair composition. Theorem~5 (Directional Asymmetry) gives an asymmetry in USD-denominated position value under symmetric sqrt-price displacement, for volatile/stablecoin pairs with the volatile asset at $T_0$. Theorem~6 (Exit Asymmetry) gives an asymmetry in USD-denominated position value between above-range and below-range exits at symmetric displacement past the boundary, for the same pair ordering. Residual magnitudes scale with position value (Theorem~1 of Ryan\cite{ryan2026a}), so the USD-denominated $|\Delta R|$ inherits the asymmetries of Theorems~5 and~6 by a scaling argument under the same domain assumptions; the bridge is developed in the remark following Theorem~5. The three results together identify when USD-denominated dust accounting acquires a directional bias even though the residual-generating mechanism is directionally blind.

The empirical confirmations below draw on Phase~3 operational data from a live CL portfolio management system on Aerodrome (Base): 36{,}337 rebalance events on 58 positions across five depositor groups over 18 days from 6--24 March 2026, executed by two contract versions (PM V8 and PM V9) on Aerodrome Slipstream.

\section{Related work}

Adams \emph{et al.}\cite{adams2021} introduced concentrated liquidity in Uniswap V3. The amount equations from that paper are the source of all three theorems below. Ryan\cite{ryan2026a} established existence, equilibrium and extinction behaviour for the geometric residual under shared depositor balances; the present paper extends that framework to directional properties.

Hasbrouck \emph{et al.}\cite{hasbrouck2025} model V3 CL provision as a covered call sold at intrinsic rather than market value, with LPs forgoing the option's time premium in exchange for fees. Theorems~5 and~6 below identify a related asymmetry. Position value in volatile/stablecoin pairs is structurally asymmetric under symmetric sqrt-price displacement, and the asymmetry intensifies at range exits. The proofs use only the V3 amount equations and the chosen num\'eraire.

Cartea \emph{et al.}\cite{cartea2024} introduce \emph{predictable loss}, a closed-form measure of the unhedgeable losses an LP incurs from holding assets in a CL pool, and analyse \emph{concentration risk}, the decrease in fee revenue when the marginal exchange rate exits the LP's range of liquidity. Theorem~6 below addresses the asymmetry of position value at the moment of range exit.

Heimbach \emph{et al.}\cite{heimbach2022} model and empirically analyse V3 LP risks and returns, finding wide cross-strategy variation. The 36{,}337-event Phase~3 dataset below adds a single-operator observation point to that record. Fan \emph{et al.}\cite{fan2023} formulate the strategic LP problem of choosing tick ranges to maximise fees net of repositioning costs. Theorems~5 and~6 describe asymmetries any such strategy faces under directional price moves and at range exits.

\section{Preliminaries}

For a position with liquidity $L$ in tick range $[s_a, s_b]$ at current sqrt price $s$, the Uniswap V3 amount equations give
\begin{equation}\label{eq:amounts-prelim}
x = L \cdot \left(\frac{1}{s} - \frac{1}{s_b}\right), \quad y = L \cdot (s - s_a)
\end{equation}
where $x$ is token0 (lower address) and $y$ is token1 (higher address). The position value in token1 units is
\begin{equation}\label{eq:value-prelim}
V = L \cdot \phi(s, s_a, s_b), \qquad \phi(s, s_a, s_b) = 2s - s^2/s_b - s_a.
\end{equation}

Theorem~1 of Ryan\cite{ryan2026a} establishes that a rebalance from range $[s_a, s_b]$ to $[s_{a'}, s_{b'}]$ at current sqrt price $s$ produces a zero geometric residual if and only if the old and new ranges demand the same token ratio at $s$. Range equality is a sufficient condition but not the only one; any range change that shifts the required token ratio produces a non-zero residual, with the sign determined by which token is the binding constraint in the new range's mint. The results below characterise how the magnitude and USD valuation of that residual respond to displacement direction, displacement magnitude and exit type.

The notation throughout is that of Ryan\cite{ryan2026a}: $s = \sqrt{p}$ is the current sqrt price; $[s_a, s_b]$ are the position range boundaries in sqrt-price space; $L$ is position liquidity; $w = s_b - s_a$ is the range width; $\delta$ is the displacement from the range centre; $nd = \delta/(w/2)$ is the normalised displacement ($0$ at the centre, $1$ at the range edge); $D$ is the dust flow (USD-denominated net position value change from non-swap sources); and $\Delta R$ is the geometric residual (token surplus from a token ratio mismatch between old and new ranges at the current price). Uniswap V3 places token0 at the pool's lower contract address. Theorems~5 and~6 below are stated for pools in which the volatile asset occupies $T_0$ (and the stablecoin $T_1$); in pools with the opposite address ordering, the USD-denominated inequalities invert but the mechanism is unchanged.

\section{Theoretical results}

\subsection{Theorem 4 (Residual Monotonicity)}

\begin{theorem}[Residual Monotonicity]
For a CL position in range $[s_a, s_b]$ rebalanced to a new range of equal width centred on the current sqrt price $s$, where $s$ has been displaced by $\delta$ from the midpoint $\bar{s} = (s_a + s_b)/2$, the absolute geometric residual $|\Delta R|$ is strictly increasing in $|\delta|$ on $(0, w/2]$, with $\Delta R(0) = 0$ and $\Delta R(w/2) = L \cdot w$. The boundary magnitude on the downward branch is smaller by a factor of $\bar{s}/s_b$ (Theorem~5).
\end{theorem}

\begin{proof}
Consider a position with liquidity $L$ in range $[s_a, s_b]$ at current sqrt price $s = \bar{s} + \delta$, where $\bar{s} = (s_a + s_b)/2$ is the range midpoint and $w = s_b - s_a$. The withdrawn token amounts are
\begin{equation}
x_w = L \cdot \left(\frac{1}{s} - \frac{1}{s_b}\right), \quad y_w = L \cdot (s - s_a)
\end{equation}

The new range, centred on $s$ with the same width $w$, has boundaries $s_{a'} = s - w/2$ and $s_{b'} = s + w/2$. At the range midpoint ($\delta = 0$), $R_{\text{old}} = R_{\text{new}}$ and the residual is zero by Theorem~1 of Ryan\cite{ryan2026a}.

For upward displacement ($\delta > 0$), $x_w$ is strictly decreasing in $\delta$ (since $\partial x_w / \partial s = -L/s^2 < 0$) whilst $y_w$ is strictly increasing ($\partial y_w / \partial s = L > 0$). As $s$ moves above the midpoint, the position becomes increasingly token1-heavy. The new range, centred on $s$, demands a ratio closer to 50/50 at its midpoint, requiring more token0 than the old position holds. Token0 is therefore the binding constraint.

The new liquidity is set by the token0 constraint:
\begin{equation}
L_{\text{new}} = \frac{x_w}{1/s - 1/s_{b'}} = \frac{L(1/s - 1/s_b)}{1/s - 1/(s + w/2)}
\end{equation}

The token1 surplus (residual) is
\begin{equation}\label{eq:dr-closed-form}
\Delta R(\delta) \;=\; y_w - L_{\text{new}}\,(s - s_{a'}) \;=\; L\,(s - s_a) \;-\; L\!\left(\frac{1/s - 1/s_b}{1/s - 1/(s + w/2)}\right)\!(w/2),
\end{equation}
with $s = \bar{s} + \delta$, $s_a = \bar{s} - w/2$, $s_b = \bar{s} + w/2$. Both subtrahend factors are strictly positive and continuously differentiable in $\delta$ on $[0, w/2)$; the leading term $L(s - s_a) = L(w/2 + \delta)$ is strictly increasing in $\delta$. A direct calculation (Appendix~B) shows that $d\,\Delta R/d\delta > 0$ on $(0, w/2)$, with $\Delta R(0) = 0$ and $\Delta R(w/2) = L \cdot w$ (the full token1 balance, the maximum possible value, at which $s = s_b$ and the position is 100\% token1).

The downward displacement case ($\delta < 0$) is symmetric with the roles of token0 and token1 exchanged: $y_w$ becomes the binding constraint, $L_{\text{new}}$ is set by token1, and the analogous closed form for $\Delta R(\delta)$ on $[-w/2, 0]$ has $d\,\Delta R/d\delta < 0$ by the same calculation. Hence $|\Delta R|$ is strictly increasing in $|\delta|$ on $(0, w/2]$.
\end{proof}

\noindent\textit{Foundry verification.} Five in-domain displacement levels on a $1{,}000$-tick-wide position at price ${\sim}\$2{,}500$ (so $|\delta| \leq w/2 = 500$ ticks).

\begin{table}[H]
\centering
\caption{Residual monotonicity, Foundry verification.}\label{tab:t4-foundry}
\begin{tabular}{@{}lrl@{}}
\toprule
Displacement (ticks) & Dust (USD) & Monotonic \\
\midrule
$+100$ & \$1{,}822 & --- \\
$+200$ & \$3{,}663 & \checkmark \\
$+300$ & \$5{,}522 & \checkmark \\
$+400$ & \$7{,}400 & \checkmark \\
$+500$ & \$9{,}297 (range exit) & \checkmark \\
\bottomrule
\end{tabular}
\end{table}

\noindent Zero violations occur across the five rows within the theorem's domain ($|\delta| \leq 500$ ticks for this 1,000-tick range). At $|\delta| = w/2$ the position fully exits range, $x_w = 0$, and the residual equals the full token1 balance ($L \cdot w$), as derived above; additional Foundry runs at $|\delta| > w/2$ (not shown) plateau at this boundary value, consistent with saturation outside the theorem's domain.

\subsection{Theorem 5 (Directional Asymmetry)}

\begin{theorem}[Directional Asymmetry]
For a CL position in pool $(T_0, T_1)$ where $T_0$ is the volatile asset and $T_1$ is a stablecoin, symmetric sqrt-price displacements $\pm \delta$ from any in-range price $s$ produce asymmetric USD-valued outcomes. The up-move position retains strictly more USD value than the down-move position at all displacement magnitudes for which both $s + \delta$ and $s - \delta$ remain within the position's range.
\end{theorem}

\begin{proof}
The CL value function in token1 (stablecoin) units \cite{ryan2026a} is
\begin{equation}
V = L \cdot \phi(s, s_a, s_b)
\end{equation}
where $\phi(s, s_a, s_b) = 2s - s^2/s_b - s_a$. When token1 is a stablecoin, $V$ is already denominated in USD.

Let $s$ be any sqrt price in the interior of the position's range. For $\delta > 0$ such that $s_a \leq s - \delta$ and $s + \delta \leq s_b$ (both displaced prices remain in range), evaluate the difference:
\begin{equation}
V(s + \delta) - V(s - \delta) = L\left[\phi(s + \delta) - \phi(s - \delta)\right]
\end{equation}

Expanding:
\begin{align}
\phi(s + \delta) - \phi(s - \delta) &= 2(s + \delta) - \frac{(s + \delta)^2}{s_b} - s_a - 2(s - \delta) + \frac{(s - \delta)^2}{s_b} + s_a \notag \\
&= 4\delta - \frac{4s\delta}{s_b} = 4\delta\left(1 - \frac{s}{s_b}\right)
\end{align}

Since $s < s_b$ strictly for any position whose current price lies in the interior of its range, the factor $(1 - s/s_b) > 0$, and therefore $V(s + \delta) > V(s - \delta)$ for all $\delta > 0$ satisfying the range constraint.
\end{proof}

\noindent\textit{Remark (Residual--value bridge).} Theorem~1 of Ryan\cite{ryan2026a} gives the residual fraction $|\Delta R|/V$ as a function of range geometry and the current price, independent of $L$; at equal $|\delta|/w$, absolute residual magnitude scales with position value. Combined with Theorem~5, this implies that USD-denominated residual magnitudes inherit the directional asymmetry proved above, under the same pair-ordering and stablecoin assumptions. Theorem~5 itself establishes only the $V$ inequality; the corresponding $|\Delta R|$ statement follows by scaling from Ryan\cite{ryan2026a}.

The asymmetry is numéraire-dependent. The proof measures value in token1 units, in which up-moves retain more value. Measured in token0 units, down-moves retain more value. The mechanism is directionally blind; the apparent bias comes from the choice of unit of account. Since USD aligns with stablecoins (typically token1 in Uniswap V3 pair ordering), USD-denominated accounting favours up-moves.

\medskip
\noindent\textit{Foundry verification.} Five symmetric displacement levels on a \$5{,}000 position; the first four rows have both displaced prices strictly inside the range (Theorem~5's domain), and the fifth ($\pm 450$ ticks) places the down-move at the lower range boundary $s_a$.

\begin{table}[H]
\centering
\caption{Directional asymmetry, Foundry verification.}\label{tab:t5-foundry}
\begin{tabular}{@{}lrrr@{}}
\toprule
Displacement & Up value (USD) & Down value (USD) & Gap \\
\midrule
$\pm 100$ ticks & \$5,522 & \$3,663 & $+$\$1,859 \\
$\pm 200$ ticks & \$6,459 & \$2,740 & $+$\$3,718 \\
$\pm 300$ ticks & \$7,400 & \$1,822 & $+$\$5,577 \\
$\pm 400$ ticks & \$8,346 & \$908 & $+$\$7,437 \\
$\pm 450$ ticks & \$8,346 & \$0 & $+$\$8,346 \\
\bottomrule
\end{tabular}
\end{table}

\noindent Zero violations across all five rows. The gap widens monotonically with displacement. The first four rows verify Theorem~5 strictly within its domain; the fifth row ($\pm 450$ ticks) is the boundary limit at which the down-move sits at $s_a$ and displays as \$0 under integer arithmetic at depressed token0 prices, giving the closing case of the inequality by continuity.

\medskip
\noindent\textit{Empirical confirmation} ($N = 26{,}191$). Events classified by normalised displacement direction.

\begin{table}[H]
\centering
\caption{Directional asymmetry, empirical confirmation.}\label{tab:t5-empirical}
\begin{tabular}{@{}lrrrr@{}}
\toprule
Direction & Events & Total Dust & Avg Dust & \% Positive \\
\midrule
Rising ($nd > 0.1$) & 17,879 & $+$\$1,129,420 & $+$\$63.17 & 35.6\% \\
Falling ($nd < -0.1$) & 8,312 & $-$\$103,308 & $-$\$12.43 & 36.2\% \\
\bottomrule
\end{tabular}
\end{table}

\noindent The proportion of events producing positive dust is near-identical (${\sim}36\%$ in both directions), so the asymmetry is in magnitude, not frequency. USD-denominated dust flow is about $5\times$ larger during rising markets, consistent with the predicted direction (token surpluses produced during rallies are valued at higher prices). The dataset is not segmented by pool type (\S\ref{sec:limitations}); the figures support the prediction at production scale without strictly confirming Theorem~5 in its volatile/stablecoin domain.

\subsection{Theorem 6 (Exit Asymmetry)}

\begin{theorem}[Exit Asymmetry]
For a CL position in pool $(T_0, T_1)$ where $T_0$ is the volatile asset and $T_1$ is a stablecoin, a below-range exit retains strictly less USD value than an above-range exit at any sqrt-price displacement $d > 0$ past the corresponding boundary.
\end{theorem}

\begin{proof}
Let $d > 0$ denote the sqrt-price displacement past the relevant range boundary, with $d < s_a$ assumed throughout so that $s' > 0$ in the below-exit case. When the current sqrt price $s'$ exits the position's range $[s_a, s_b]$, the V3 amount equations reduce to boundary values:

At below-exit ($s' = s_a - d$, position is 100\% token0):
\begin{equation}
x_{\text{below}} = L \cdot (1/s_a - 1/s_b), \quad y_{\text{below}} = 0
\end{equation}

At above-exit ($s' = s_b + d$, position is 100\% token1):
\begin{equation}
x_{\text{above}} = 0, \quad y_{\text{above}} = L \cdot (s_b - s_a)
\end{equation}

The above-exit position holds only token1, so $V_{\text{above}} = y_{\text{above}} = L(s_b - s_a)$ in token1 units; this equals USD since token1 is a stablecoin, and the value is constant in $d$.

The USD value of the below-exit position is $V_{\text{below}} = x_{\text{below}} \cdot s'^2 = L(1/s_a - 1/s_b)(s_a - d)^2$, where $s'^2$ is the USD price of token0 at the current market sqrt price. This is strictly decreasing in $d$ on $(0, s_a)$.

For any $d > 0$, the ratio $V_{\text{below}}/V_{\text{above}}$ decreases with $d$. The below-exit position loses value as the price moves further past the boundary, whilst the above-exit position does not.
\end{proof}

\noindent\textit{Remark (Redeployment).} To redeploy either terminal position into a new in-range position, the missing token must be obtained via swap. The below-exit position holds only token0, which has depreciated; the above-exit position holds only token1, which retains its full stablecoin value. Since $V_{\text{below}}$ decreases with $d$ whilst $V_{\text{above}}$ is constant, the below-exit position has strictly less capital available for redeployment at any $d > 0$.

\noindent\textit{Foundry verification.} Five symmetric exit distances past the range boundary.

\begin{table}[H]
\centering
\caption{Exit asymmetry, Foundry verification.}\label{tab:t6-foundry}
\begin{tabular}{@{}lrrrr@{}}
\toprule
Exit distance & Above value & Below value & Above swap frac & Below swap frac \\
\midrule
$\pm 50$ ticks & \$9,297 & \$0 & 0 bps & 9,999 bps \\
$\pm 100$ ticks & \$9,297 & \$0 & 0 bps & 9,999 bps \\
$\pm 200$ ticks & \$9,297 & \$0 & 0 bps & 9,999 bps \\
$\pm 300$ ticks & \$9,297 & \$0 & 0 bps & 9,999 bps \\
$\pm 500$ ticks & \$9,297 & \$0 & 0 bps & 9,999 bps \\
\bottomrule
\end{tabular}
\end{table}

\noindent The above-exit position retains its full stablecoin value (\$9{,}297) at all five distances. The below-exit position displays as \$0 in the test output because its USD value falls below the one-cent display threshold under integer arithmetic at depressed token0 prices; the token0 quantity is non-zero but negligible relative to the above-exit value. The swap fraction columns report how much of the held token must be converted to obtain the missing token for redeployment: above-swap = 0~bps (the above-exit position already holds the stablecoin), below-swap = 9{,}999~bps (the below-exit position must convert almost everything). Both $V_{\text{above}} > V_{\text{below}}$ and the swap fraction gap hold at all five distances (ten verifications). The strict-monotonicity-in-$d$ claim of the proof is masked here by the integer-arithmetic floor on $V_{\text{below}}$; what the table verifies is the qualitative inequality $V_{\text{above}} > V_{\text{below}}$ at every $d > 0$, not the rate of decay.

\medskip
\noindent\textit{Empirical confirmation} ($N = 2{,}225$ exit events).

\begin{table}[H]
\centering
\caption{Exit asymmetry, empirical confirmation.}\label{tab:t6-empirical}
\begin{tabular}{@{}lrrrrr@{}}
\toprule
Exit type & Events & Avg Dust & Median & P5 (worst) & P95 (best) \\
\midrule
Above & 1,305 & $-$\$89 & $-$\$26 & $-$\$1,261 & $+$\$609 \\
Below & 920 & $-$\$141 & $-$\$29 & $-$\$1,479 & $+$\$961 \\
\bottomrule
\end{tabular}
\end{table}

\noindent In this dataset, below-exits produce $58\%$ larger average negative dust flow. The distribution tails are wider in both directions; P5 for below-exits is $-$\$1{,}479 versus $-$\$1{,}261 for above ($17\%$ wider). Average tick displacement at exit is $+103$ ticks (above) versus $-88$ ticks (below). Theorem~6 establishes $V_{\text{below}} < V_{\text{above}}$ at the terminal state within its volatile/stablecoin domain. The observed dust flow asymmetry is directionally consistent with this prediction via the residual--value bridge (Remark after Theorem~5). Theorem~6 does not prove a direct $|\Delta R|$ asymmetry for exit rebalances, and the dataset is unsegmented by pool composition (\S\ref{sec:limitations}).

\section{Discussion}

Theorem~4 gives strict monotonicity of $|\Delta R|$ in displacement up to the range boundary. Within their volatile/stablecoin domain, Theorems~5 and~6 show that position value is directionally asymmetric in USD terms even though the residual-generating mechanism is directionally blind. Residual magnitudes scale with position value, so USD-denominated residuals inherit these asymmetries (remark after Theorem~5).

The siphon moves tokens between positions through a shared dust pool based on geometric ratios. It does not take directional views. When the token flows are valued in USD, the accounting reflects the price exposure of the tokens being moved. In a portfolio where a volatile asset appears in multiple pools paired against stablecoins (a hub topology of the kind Theorems~5 and~6 cover), token surpluses generated by rebalancing carry that asset's price exposure. The USD value of those flows is predicted to rise during rallies relative to selloffs at comparable displacement (Theorem~5 via the residual--value bridge), and below-range exits are predicted to retain less USD value than above-range exits at symmetric displacement past the boundary (Theorem~6). For portfolios with this composition, USD-denominated performance metrics reflect the hub token's directional exposure. In pools with the opposite address ordering or non-stablecoin pairings, the asymmetry inverts or attenuates with the numéraire.

\section{Foundry verification suite}\label{sec:foundry}

The directional theorems were verified against a mock pool implementing exact Uniswap V3 TickMath constants for sqrt-price and liquidity computations (\texttt{MockCLPoolV2}). The full Foundry suite is shared with the companion paper \cite{ryan2026a} and runs as 15 tests across four contracts:

\begin{codeblock}
\begin{verbatim}
cd foundry
forge test -vv --no-match-contract 'GeometricResidualProof$'
Suite result: ok. 10 passed; 0 failed; 0 skipped     (mock-pool tests)

RPC_BASE_ALCHEMY=https://mainnet.base.org forge test -vv
Suite result: ok. 5 passed; 0 failed; 0 skipped      (live fork tests)
\end{verbatim}
\end{codeblock}

The mock-pool subset (10 tests) verifies existence and the scaling-with-position-size property of the residual (Theorem~1), zero-swap extinction (Theorem~3), and the three directional theorems below. The live fork subset (5 tests) reproduces the Theorem~1 scenarios against unmodified Aerodrome Slipstream contracts on Base mainnet via \texttt{vm.createSelectFork}, exercising the same withdraw-mint sequence on real on-chain V3 geometry. The qualitative claims (residual existence, same-range control, residual scaling with position size at fixed range geometry) hold deterministically; the exact residual values vary between runs as live pool state changes.

Tests 8--10 below verify the directional theorems introduced in this paper:

\begin{table}[H]
\centering
\caption{Tests for the directional theorems.}\label{tab:test-summary}
\begin{tabular}{@{}llll@{}}
\toprule
Test & Theorem & Assertion & Result \\
\midrule
8  & Theorem 4 (monotonicity)          & 0 violations, 5 in-domain levels & PASSED \\
9  & Theorem 5 (directional asymmetry) & UP $>$ DOWN at 4 in-domain + 1 boundary & PASSED \\
10 & Theorem 6 (exit asymmetry)        & 10/10 confirmations          & PASSED \\
\bottomrule
\end{tabular}
\end{table}

\noindent Tests 1--7 (Theorems 1 and 3 from Ryan\cite{ryan2026a}) and the five live fork tests are included for completeness and all pass.

\section{Limitations}\label{sec:limitations}

Four limitations should be flagged. \emph{Empirical scope.} All empirical data come from one operator's positions on Aerodrome (Base). The theorem proofs are general, derived from the V3 amount equations, but the empirical ratios ($5\times$ directional asymmetry, 58\% exit asymmetry) are specific to this dataset and may differ under other protocols, pairs or market conditions.

\emph{Observation window.} 18 days of continuous operation. The dataset does not span a full market cycle.

\emph{Unsegmented pool composition.} The empirical confirmations for Theorems~5 and~6 ($N = 26{,}191$ directional events and $N = 2{,}225$ exit events) are drawn from the full Phase~3 dataset without segmentation by pool type. The reported ratios may therefore conflate volatile/stablecoin pools, to which the theorems strictly apply, with any other pool compositions present. Segmented analysis is left to future work.

\emph{Num\'eraire dependence.} Theorems~5 and~6 are proved in token1 units. Measured in token0 units, the inequalities reverse. The directional asymmetry is therefore an artefact of the chosen unit of account, not an intrinsic property of the mechanism. The USD interpretation holds when token1 is a stablecoin or when token1's USD price is approximately stable over the rebalance interval.

\section{Conclusion}

In concentrated liquidity rebalancing, the geometric residual increases strictly monotonically with displacement up to the range boundary (Theorem~4, Residual Monotonicity). For volatile/stablecoin pairs with the volatile asset at $T_0$, position value is directionally asymmetric in USD terms under symmetric in-range displacement (Theorem~5, Directional Asymmetry); USD-denominated residuals inherit this asymmetry by a scaling argument from the residual fraction result of Ryan\cite{ryan2026a}. For the same pair ordering, below-range exits retain less USD value than above-range exits at symmetric displacement past the boundary (Theorem~6, Exit Asymmetry). All three results follow from the Uniswap V3 amount equations and are verified by Foundry tests using exact V3 \texttt{TickMath} constants for sqrt-price and liquidity computations.

The empirical dataset of 36{,}337 events across 58 positions and 18 days is consistent with the predicted directions at production scale. USD-denominated dust flows are about $5\times$ larger in rising markets than in falling markets despite a near-identical frequency of positive dust events, and below-range exits show $58\%$ larger average negative dust flow than above-range exits at comparable displacement. The dataset is not segmented by pool composition (\S\ref{sec:limitations}); the figures support the predictions at production scale without strictly confirming Theorems~5 and~6 in their volatile/stablecoin domain.

Ryan\cite{ryan2026a} established existence (Theorem~1), equilibrium (Theorem~2) and extinction (Theorem~3) of the geometric residual under shared depositor balances. The three theorems above extend that framework to displacement magnitude, price direction and exit type. The proofs follow from the V3 amount equations and the chosen num\'eraire; under the stated domain assumptions, the same properties should manifest in any system performing autonomous CL rebalancing with shared depositor balances.

\addcontentsline{toc}{section}{References}
\begin{thebibliography}{9}

\bibitem{ryan2026a}
Ryan, K.R. ``The Geometric Siphon: Emergent Capital Reallocation in Concentrated Liquidity Portfolios.'' \emph{SSRN preprint 6374838} (2026) \url{https://doi.org/10.2139/ssrn.6374838}.

\bibitem{adams2021}
Adams, H., Zinsmeister, N., Salem, M., Keefer, R., Robinson, D. ``Uniswap v3 Core.'' \emph{Uniswap Labs Technical Whitepaper} (2021) \url{https://uniswap.org/whitepaper-v3.pdf}.

\bibitem{hasbrouck2025}
Hasbrouck, J., Rivera, T.J., Saleh, F. ``An Economic Model of a Decentralized Exchange with Concentrated Liquidity.'' \emph{Management Science (forthcoming)} (2025) \url{https://doi.org/10.1287/mnsc.2024.04510}.

\bibitem{cartea2024}
Cartea, \'A., Drissi, F., Monga, M. ``Decentralized Finance and Automated Market Making: Predictable Loss and Optimal Liquidity Provision.'' \emph{SIAM Journal on Financial Mathematics (forthcoming)} (2024) \url{https://doi.org/10.1137/23M1602103}.

\bibitem{heimbach2022}
Heimbach, L., Schertenleib, E., Wattenhofer, R. ``Risks and Returns of Uniswap V3 Liquidity Providers.'' \emph{Proceedings of the 4th ACM Conference on Advances in Financial Technologies (AFT 2022)} (2022) \url{https://doi.org/10.1145/3558535.3559772}.

\bibitem{fan2023}
Fan, Z., Marmolejo-Cossío, F.J., Moroz, D., Neuder, M., Rao, R., Parkes, D.C. ``Strategic Liquidity Provision in Uniswap v3.'' \emph{Proceedings of the 5th Conference on Advances in Financial Technologies (AFT 2023), LIPIcs vol.~282, Article 25} (2023) \url{https://doi.org/10.4230/LIPIcs.AFT.2023.25}.

\end{thebibliography}

\appendix
\titleformat{\section}[block]{\fontsize{13pt}{15pt}\bfseries}{Appendix \thesection:}{0.6em}{}
\titlespacing*{\section}{0pt}{18pt}{8pt}

\section{Foundry test reproduction}\label{app:foundry}

Requirements: Foundry \texttt{forge} (installation instructions at \url{https://book.getfoundry.sh}) and, for the fork tests, any working Base RPC URL. The \texttt{forge-std} library is included as a git submodule; the \texttt{git submodule update} command below fetches it into \texttt{foundry/lib/forge-std}.

\begin{codeblock}
\begin{verbatim}
cd foundry
git submodule update --init --recursive

# Mock-pool tests only, no network access required (10 tests)
forge test -vv \
  --no-match-contract '(GeometricResidualProof|DirectionalExitForkProof)$'

# Full suite, including the 5 live fork tests (15 tests)
RPC_BASE_ALCHEMY=https://mainnet.base.org forge test -vv

# Directional theorem tests only (Tests 8-10 in NewTheoremsProof.t.sol)
forge test -vv --match-contract NewTheoremsProof
# Test 8: Monotonicity            0 violations across 5 in-domain levels
# Test 9: Directional asymmetry   5 symmetric moves, UP > DOWN at all levels
# Test 10: Exit asymmetry         10/10 below < above

# MockCLPoolV2.sol uses exact Uniswap V3 TickMath constants.
\end{verbatim}
\end{codeblock}

\noindent The official public endpoint above is sufficient for all five fork tests; any Alchemy, Infura, QuickNode or other Base RPC URL is equally acceptable in \texttt{RPC\_BASE\_ALCHEMY}. Forge caches RPC responses locally, so repeat runs are fast; mock-pool tests complete in under a second. The captured \texttt{forge test -vv} output mapped to each theorem lives in \texttt{foundry/PROOF\_OUTPUT.md} and serves as the deterministic regression target for the mock-pool tests; fork-test residuals vary between runs as live pool state changes (the qualitative claims hold deterministically).

\section[Derivation for Theorem 4]{Derivation of $d\,\Delta R/d\delta$ for Theorem~4}\label{app:dr-derivation}

The closed form for the token1 surplus on $\delta \in [0, w/2]$, from equation~\eqref{eq:dr-closed-form}, is
\begin{equation}
\Delta R(\delta) \;=\; L\,(s - s_a) \;-\; L\!\left(\frac{1/s - 1/s_b}{1/s - 1/(s + w/2)}\right)\!(w/2),
\end{equation}
with $s = \bar{s} + \delta$, $s_a = \bar{s} - w/2$, $s_b = \bar{s} + w/2$, $w = s_b - s_a$. Define
\begin{equation}
g(\delta) \;\equiv\; \frac{1/s - 1/s_b}{1/s - 1/(s + w/2)}.
\end{equation}
Using $1/s - 1/s_b = (s_b - s)/(s\,s_b)$ and $1/s - 1/(s + w/2) = (w/2)/(s\,(s + w/2))$, the factors of $s$ cancel and
\begin{equation}\label{eq:g-simplified}
g(\delta) \;=\; \frac{(s_b - s)\,(s + w/2)}{(w/2)\,s_b} \;=\; \frac{(w/2 - \delta)\,(\bar{s} + \delta + w/2)}{(w/2)\,(\bar{s} + w/2)}.
\end{equation}
Substituting $s = \bar{s} + \delta$ and $s_b = \bar{s} + w/2$ gives $\Delta R$ as an explicit rational function of $\delta$:
\begin{equation}
\Delta R(\delta) \;=\; L\,(w/2 + \delta) \;-\; L\,\frac{(w/2 - \delta)\,(\bar{s} + \delta + w/2)}{\bar{s} + w/2}.
\end{equation}
Applying the product rule to the second term, with $A(\delta) = w/2 - \delta$ and $B(\delta) = \bar{s} + \delta + w/2$ so that $A'(\delta) = -1$ and $B'(\delta) = 1$:
\begin{equation}
\frac{d}{d\delta}\bigl[A(\delta)\,B(\delta)\bigr] \;=\; -B(\delta) + A(\delta) \;=\; -(\bar{s} + \delta + w/2) + (w/2 - \delta).
\end{equation}
Therefore
\begin{equation}
\frac{d\,\Delta R}{d\delta} \;=\; L \;-\; \frac{L\bigl[\,-B(\delta) + A(\delta)\,\bigr]}{\bar{s} + w/2} \;=\; \frac{L\bigl[(\bar{s} + w/2) + B(\delta) - A(\delta)\bigr]}{\bar{s} + w/2}.
\end{equation}
Substituting and collecting,
\begin{equation}
(\bar{s} + w/2) + (\bar{s} + \delta + w/2) - (w/2 - \delta) \;=\; 2\bar{s} + 2\delta + w/2,
\end{equation}
so
\begin{equation}
\frac{d\,\Delta R}{d\delta} \;=\; L\,\frac{2\bar{s} + 2\delta + w/2}{\bar{s} + w/2}.
\end{equation}
For any economically meaningful CL position, $\bar{s} > 0$, $w > 0$ and $\delta \in (0, w/2]$, so $2\bar{s} + 2\delta + w/2 > 0$ and $\bar{s} + w/2 > 0$. Hence
\begin{equation}
\frac{d\,\Delta R}{d\delta} \;>\; 0 \quad\text{on}\quad (0, w/2],
\end{equation}
and $\Delta R$ is strictly increasing on $[0, w/2]$. At the endpoints, $\Delta R(0) = L(w/2) - L(w/2)(\bar{s} + w/2)/(\bar{s} + w/2) = 0$, consistent with Theorem~1 of Ryan\cite{ryan2026a}; and $\Delta R(w/2) = L \cdot w$, the full token1 balance at range exit.

The downward case $\delta \in [-w/2, 0]$ proceeds analogously with token1 as the binding constraint: $L_{\text{new}}$ is set by $y_w$ rather than $x_w$, $g(\delta)$ takes the symmetric form, and the same calculation gives $d\,\Delta R/d\delta < 0$ on $[-w/2, 0)$. Combining the two cases, $|\Delta R(\delta)|$ is strictly increasing in $|\delta|$ on $(0, w/2]$, as claimed in Theorem~4.

\section{Data and reproduction}\label{app:reproduction}

The repository is at \url{https://github.com/g4nc0r/the-geometric-siphon}; the Zenodo Code DOI \url{https://doi.org/10.5281/zenodo.19905208} resolves to an immutable archival snapshot, whilst the GitHub URL tracks the latest commit. The current Foundry suite and reproduction harness in the repository align with the consolidated manuscript referenced in the supersession note. The qualitative theorem verifications described in \S\ref{sec:foundry} (residual monotonicity in displacement, directional asymmetry of USD-denominated position value under symmetric in-range moves, and exit asymmetry between above-range and below-range exits) reproduce on the current harness; specific numerical values may differ from the figures reported here, which reflect the harness used at the time of this preprint. The Phase~3 diffusion-event log of 36{,}337 rebalance events across 58 positions used for the empirical confirmations is also deposited on Zenodo (Data DOI \url{https://doi.org/10.5281/zenodo.19904832}).

\end{document}
